%% kapa: exam=1 round=22 year=2011 period=2 form=A subject=appraisal_law subject_ko=부동산관계법규 seq=25 number=65
%% source: Q-Net 감정평가사 기출문제 (cst00309) artlSeq=5036790 file=extracted/1008/5036790/2011년 제22회 감정평가사 1차시험 기출문제/2교시.hwp
%% conversion: kordoc 4.24.51 -> tools/gen1.py (auto); review=proofread-gpt-astra; flags=none
\begin{question}{65}
국유재산법령상 국유재산의 관리에 관한 설명으로 옳지 않은 것은?
\begin{choices}
  \item 기획재정부장관은 일반재산을 보존용재산으로 전환하여 관리할 수 있다.
  \item 행정재산이라도 판결에 따라 사권을 설정하는 것은 허용된다.
  \item 판결에 따라 취득하는 경우에는 사권이 설정된 재산이라도 국유재산으로 취득할 수 있다.
  \item 국가는 행정재산에 건물, 교량 등 구조물과 그 밖의 영구시설물을 축조할 수 있다.
  \item 총괄청이나 중앙관서의 장은 소유자 없는 부동산을 국유재산으로 취득한다.
\end{choices}
\end{question}
\answer{2}
